\section{Point 2 — Relationship Between $f_s$, $P$, $S$, $f$; Phase-Locking and Degeneracies}


\subsection{The Patch-Induced Frequency}

The tokeniser introduces a characteristic \emph{patch frequency}:
\begin{equation}
    f_{\mathrm{patch}} = \frac{f_s}{S}.
\end{equation}
This is the rate at which new patches are produced (in Hz, referred to original time),
analogous to a resampling rate in the token domain.

\subsection{Phase-Locking Condition}

Consider a sinusoidal signal $x[n] = A\cos\!\bigl(2\pi f n / f_s + \phi\bigr)$.

\begin{definition}[Phase-Locking]
    The signal \emph{phase-locks} to the patch grid when its period $T_0 = f_s/f$ (in
    samples) is \emph{commensurate} with $S$:
    \begin{equation}
        \frac{S \cdot f}{f_s} \in \mathbb{Q}.
    \end{equation}
    The degenerate (critical) case occurs when this ratio is an integer:
    \begin{equation}
        f = m \cdot \frac{f_s}{S} = m \cdot f_{\mathrm{patch}},
        \qquad m \in \mathbb{Z}^+,
    \end{equation}
    i.e.\ $f$ is a \textbf{harmonic of} $f_{\mathrm{patch}}$.
\end{definition}

\subsection{Effect of Patch Size $P$}

A single patch of $P$ samples spans a duration $P/f_s$ seconds and covers
$P \cdot f / f_s$ cycles of the signal.  When
\begin{equation}
    \frac{P \cdot f}{f_s} \in \mathbb{Z}^+,
    \quad \text{i.e.} \quad
    f = k \cdot \frac{f_s}{P}, \quad k \in \mathbb{Z}^+,
\end{equation}
the patch contains an exact integer number of cycles: its content is \emph{internally
redundant}.

The most severe degeneracy occurs when both the stride and the patch-size conditions hold
simultaneously:
\begin{equation}
    f = \frac{\ell \cdot f_s}{\mathrm{lcm}(P,\, S)},
    \qquad \ell \in \mathbb{Z}^+.
\end{equation}

\subsection{Token-Space Degeneracies}

\paragraph{At $f = m\,f_s/S$ (harmonics of $f_{\mathrm{patch}}$).}
\begin{itemize}
    \item All patches are translates of the same waveform by a phase that is a multiple of
          $2\pi$: they become \textbf{identical} (cf.\ Point~1).
    \item The token sequence is constant; it encodes \textbf{no phase information}.
    \item Two signals at the same frequency but different phases (shifted by any multiple of
          $S$ samples) map to the same token sequence: \emph{token-space invariance}.
\end{itemize}

\paragraph{At $f = k\,f_s/P$ (but not harmonics of $f_{\mathrm{patch}}$).}
\begin{itemize}
    \item Each patch is internally periodic, so its content is \emph{within-patch redundant}.
    \item Inter-patch variation still exists (partial degeneracy).
\end{itemize}

\subsection{Overlap, Stride, and Blind-Spot Frequencies}

When $S < P$ (overlapping patches), consecutive patches share $P - S$ samples, partially
mitigating the aliasing effect even near harmonic frequencies.  For $S \geq P$ (no overlap
or gaps), the aliasing is maximal; the \textbf{blind-spot frequencies} are:
\begin{equation}
    \boxed{f_{\mathrm{blind}} = m \cdot \frac{f_s}{S},
    \qquad m = 1, 2, 3, \ldots, \left\lfloor \frac{S}{2} \right\rfloor.}
\end{equation}
The upper limit follows from the Nyquist criterion applied to the token sequence, whose
effective sample rate is $f_s/S$.

\subsection{Summary}

\begin{table}[h]
    \centering
    \begin{tabular}{@{}ll@{}}
        \toprule
        \textbf{Condition} & \textbf{Effect} \\
        \midrule
        $f = m\,f_s/S$ & Full phase-locking $\to$ identical patches $\to$ constant token sequence \\
        $f = k\,f_s/P$ & Internal patch periodicity $\to$ within-patch redundancy \\
        Both simultaneously & Maximal degeneracy \\
        $S < P$ (overlap) & Partial mitigation; residual phase information \\
        $S \geq P$ (no overlap) & Full degeneracy at blind-spot frequencies \\
        \bottomrule
    \end{tabular}
    \caption{Summary of conditions and their effects on token-space observability.}
\end{table}
